\title{DeepSeekMath: Pushing the Limits of Mathematical Reasoning in Open Language Models}

\begin{abstract}
The inherently intricate and systematic characteristics of mathematical reasoning present considerable difficulties for language models. In this study, we present DeepSeekMath~7B, a model that extends DeepSeek-Coder-Base-v1.5 7B by engaging in further pre-training on 120B math-focused tokens, collected from Common Crawl as well as natural language and code sources. DeepSeekMath~7B attains a notable 51.7\% accuracy on the MATH benchmark at the competition level without the assistance of external toolkits or ensemble methods, narrowing the gap to models such as Gemini-Ultra and GPT-4. Additionally, DeepSeekMath~7B reaches 60.9\% with self-consistency across 64 generated samples on the same benchmark. The superior mathematical reasoning of DeepSeekMath~can be traced to two principal elements: first, the effective exploitation of public web datasets via a precisely designed pipeline for data selection; second, the introduction of Group Relative Policy Optimization (GRPO)—a modified approach to Proximal Policy Optimization (PPO)—which both bolsters mathematical problem-solving capabilities and streamlines PPO memory consumption.
\end{abstract}

\section{Introduction}

Large language models (LLMs) have transformed artificial intelligence's capacity for mathematical reasoning, leading to substantial progress on quantitative reasoning benchmarks \citep{MATH} as well as on geometry-specific assessments \citep{trinh2024solving}. Additionally, LLMs have demonstrated their value as tools for aiding humans in tackling intricate mathematical tasks \citep{tao}. Despite these developments, state-of-the-art systems like GPT-4 \citep{gpt4} and Gemini-Ultra \citep{gemini} remain inaccessible to the public, while existing open-source alternatives fall considerably short in terms of performance.

This paper presents DeepSeekMath, a specialized language model designed for mathematical reasoning that markedly surpasses the capabilities of current open-source models and nearly matches GPT-4's performance on academic mathematics tasks. Our method involves constructing the DeepSeekMath Corpus, a comprehensive, high-quality dataset containing 120B tokens of mathematical content, curated from the Common Crawl (CC) via a fastText classifier \citep{joulin2016fasttext}. Initially, the classifier is built with examples from OpenWebMath \citep{openwebmath} as positive samples and a diverse range of other websites as negative samples. We then utilize this classifier to extract additional positive mathematical data from CC, which are further vetted by human annotators. The classifier is iteratively improved using this expanded and curated dataset. Our experimental findings demonstrate that the resulting large-scale corpus exhibits strong quality: the base model DeepSeekMath-Base 7B achieves 64.2\% accuracy on the GSM8K benchmark \citep{gsm8k} and 36.2\% on the MATH competition-level dataset \citep{MATH}, exceeding the performance of Minerva 540B \citep{minerva}. Furthermore, as the DeepSeekMath Corpus supports multiple languages, we also observe notable gains on Chinese mathematical evaluations \citep{wei2023cmath,agieval}. We propose that our approach to mathematical data collection and refinement can serve as a useful foundation for further work in the field, and we anticipate substantial opportunities for future enhancement.

DeepSeekMath-Base utilizes DeepSeek-Coder-Base-v1.5 7B \citep{deepseek-coder} as its foundation, reflecting our observation that initializing from a code-pretrained model yields superior results over beginning with a generic large language model. Additionally, our findings reveal that training on mathematical data enhances performance on the MMLU \citep{mmlu} and BBH \citep{bbh} benchmarks, suggesting that such training contributes not only to the improvement of mathematical competence but also to the model’s overall reasoning skills.

Following the pre-training phase, we conduct mathematical instruction tuning on DeepSeekMath-Base utilizing chain-of-thought \citep{cot}, program-of-thought \citep{pot,pal}, and tool-integrated reasoning \citep{tora} datasets. The resulting DeepSeekMath-Instruct 7B model outperforms all other models in the 7B parameter range and demonstrates performance on par with instruction-tuned open-source models at the 70B scale.

Additionally, we propose Group Relative Policy Optimization (GRPO), a reinforcement learning (RL) method derived from Proximal Policy Optimization (PPO) \citep{schulman2017proximal}. Unlike conventional approaches, GRPO omits the critic model by deriving the baseline directly from group scores, leading to notable savings in training computation. Even when trained on only a fraction of English instruction tuning data, GRPO achieves marked improvements over the already robust DeepSeekMath-Instruct, showing higher performance on both in-domain (GSM8K: 82.9\% $\rightarrow$ 88.2\%, MATH: 46.8\% $\rightarrow$ 51.7\%) and out-of-domain mathematical benchmarks (for instance, CMATH: 84.6\% $\rightarrow$ 88.8\%) throughout the RL process. We further articulate a unified conceptual framework that elucidates the connections among Rejection Sampling Fine-Tuning (RFT) \citep{yuan2023scaling}, Direct Preference Optimization (DPO) \citep{dpo}, PPO, and GRPO, showing that these approaches can all be viewed as variants or simplifications of RL methodologies. To scrutinize the fundamental aspects of this framework, we conduct comprehensive analyses, contrasting, for example, online versus offline learning, outcome-based versus process-based supervision, and single-step versus iterative RL. Lastly, we offer a rationale for how RL enhances the capabilities of instruction-tuned models, and discuss prospective research avenues for advancing RL within this unified theoretical perspective.

\subsection{Contributions}
We present scalable mathematical pre-training methods, complemented by an in-depth examination and evaluation of reinforcement learning approaches.

\textbf{Math Pre-Training at Scale}

\begin{itemize}[topsep=0pt]
    \item In this study, we demonstrate that the openly available Common Crawl dataset harbors significant resources relevant to mathematical research. Through the deployment of a carefully crafted data curation process, we assemble the DeepSeekMath Corpus, a rigorously filtered and high-quality collection comprising 120B tokens drawn from web pages with mathematical relevance. This corpus is nearly sevenfold larger than the set of mathematical web pages used in Minerva \citep{minerva}, and approximately nine times greater than the recently introduced OpenWebMath \citep{openwebmath}.

\item The DeepSeekMath-Base 7B model, which has undergone pre-training, attains a level of performance similar to that of Minerva 540B \citep{minerva}. This suggests that sheer parameter count does not solely determine proficiency in mathematical reasoning; rather, a relatively compact model trained on well-curated data can also deliver robust results.

\item We present results from our experiments focused on mathematical training. Engaging in code training before introducing math tasks enhances the model's performance on mathematical problem solving, irrespective of tool use. These observations contribute to the ongoing debate around the question: \textit{does code training enhance reasoning skills?} Our results indicate that it does, specifically in the context of mathematical reasoning.

\item While the use of arXiv papers for training is widespread, particularly among numerous mathematics-focused studies, it does not result in any significant enhancements across the mathematical benchmarks utilized in this work.
\end{itemize}

\textbf{Exploration and Analysis of Reinforcement Learning}

\begin{itemize}[topsep=0pt]
    \item This paper presents Group Relative Policy Optimization (GRPO), a resource-efficient and high-performing reinforcement learning technique. Unlike traditional approaches, GRPO eliminates the need for a critic network by leveraging group-based score baselines, thereby reducing computational demands relative to Proximal Policy Optimization (PPO).
    \item Through empirical studies, we show that GRPO notably boosts the abilities of our instruction-finetuned model DeepSeekMath-Instruct, utilizing only instruction-tuning datasets. Additionally, we find that GRPO yields improvements in generalization to out-of-domain tasks during reinforcement learning.
    \item We offer a comprehensive framework that encompasses various approaches, including RFT, DPO, PPO, and GRPO. A wide array of experimental comparisons—including online versus offline learning, outcome-based versus process-based feedback, and single-step versus iterative RL—are performed to dissect the fundamental components of this unified view.
    \item Leveraging this unified perspective, we analyze the mechanisms that underpin the success of reinforcement learning, and outline several promising avenues for advancing reinforcement learning techniques for large language models (LLMs).
\end{itemize}

\subsection{Summary of Evaluations and Metrics}

\begin{itemize}[topsep=0pt]
    \item \textbf{English and Chinese Mathematical Reasoning}: Our study extensively evaluates model performance on a variety of English and Chinese mathematics benchmarks, spanning problems from elementary to collegiate levels. For English-language assessment, we include GSM8K \citep{gsm8k}, MATH \citep{MATH}, SAT \citep{llemma}, OCW Courses \citep{minerva}, and MMLU-STEM \citep{mmlu}. In the Chinese context, the benchmarks utilized are MGSM-zh \citep{mgsm}, CMATH \citep{wei2023cmath}, Gaokao-MathCloze \citep{agieval}, and Gaokao-MathQA \citep{agieval}. We test model competence both in generating detailed, self-explanatory solutions in text, independent of external tools, as well as in leveraging Python for computational problem-solving.

On English benchmark evaluations, DeepSeekMath-Base rivals the proprietary Minerva 540B model \citep{minerva}, and outperforms all open-source base models—including Mistral 7B \citep{mistral} and Llemma-34B \citep{llemma}—irrespective of their engagement in math-specific pre-training, frequently by a considerable margin. Importantly, DeepSeekMath-Base demonstrates even greater proficiency on Chinese benchmarks, a result that can be attributed to our departure from previous practices \citep{minerva,llemma} of exclusively sourcing English mathematical pre-training data; instead, our dataset also incorporates high-quality math content from non-English sources. Leveraging mathematical instruction tuning alongside reinforcement learning, our subsequent models, DeepSeekMath-Instruct and DeepSeekMath-RL, achieve remarkable results, with the latter surpassing the 50\% accuracy threshold on the competition-grade MATH dataset—setting a new milestone for open-source systems.

\item \textbf{Formal Mathematics}: We assess DeepSeekMath-Base's capabilities on the informal-to-formal theorem proving benchmark from \citep{dsp_proof}, employing miniF2F \citep{minif2f} and selecting Isabelle \citep{isabelle} as the proof assistant. DeepSeekMath-Base achieves notable few-shot performance in the autoformalization setting.

\item \textbf{Natural Language Understanding, Reasoning, and Code}: To thoroughly evaluate the breadth of the models' understanding, reasoning skills, and programming ability, we test DeepSeekMath-Base on the Massive Multitask Language Understanding (MMLU) benchmark \citep{mmlu}, which features 57 multiple-choice problems across a variety of domains. Additionally, we employ BIG-Bench Hard (BBH) \citep{bbh}, encompassing 23 complex tasks with an emphasis on multi-step reasoning, alongside HumanEval \citep{codex} and MBPP \citep{mbpp}, both standard benchmarks for code generation. Pre-training on mathematics contributes to improved performance in both linguistic understanding and reasoning tasks.

\section{Math Pre-Training}

\subsection{Data Collection and Decontamination} Here, we describe the methodology employed in building the DeepSeekMath Corpus using data sourced from Common Crawl. As illustrated in Figure \ref{fig:data_collect}, an iterative pipeline is introduced to enable the structured extraction of an extensive mathematical dataset from Common Crawl, beginning with a seed corpus comprised of a limited yet highly curated set of math-related datasets. It is important to emphasize that this strategy is versatile and can be similarly utilized for various domains, including programming.

To begin, we select OpenWebMath \citep{openwebmath}—a curated compilation of high-quality mathematical content available on the web—as our foundational seed dataset. We utilize this collection to train a fastText model \citep{joulin2016fasttext}, aiming to retrieve additional web pages with similar mathematical characteristics to those in OpenWebMath. Our approach involves sampling 500,000 entries from the seed corpus to serve as positive instances and pairing them with 500,000 randomly chosen web pages from Common Crawl as negative instances. Training is conducted using an open-source implementation\footnote{\url{https://fasttext.cc}} of fastText, where the model is configured with a 256-dimensional embedding space, a learning rate of 0.1, word n-grams up to length 3, a minimum word frequency threshold of 3, and three epochs of training. To manage the scale of Common Crawl, we implement both URL-level deduplication and near-deduplication, ultimately reducing the dataset to 40B unique HTML web pages. Employing the trained fastText model, we identify candidate mathematical pages from this pruned collection. To ensure quality, we rank these candidate pages based on their fastText model scores and retain only those with the highest rankings. To determine how much data to include, we conduct preliminary pre-training assessments across datasets comprising the top 40B, 80B, 120B, and 160B tokens. In our initial iteration, we opt to retain the dataset corresponding to the highest-scoring 40B tokens.

Following the initial round of data collection, a significant number of mathematical web pages remain uncollected. This gap is largely attributable to the limited variety within the set of positive examples used to train the fastText model. To address this, we expand the seed corpus by incorporating additional mathematical web sources in order to enhance the performance of the fastText model. Our method involves first partitioning the entire Common Crawl dataset into distinct domains, where each domain comprises web pages with a common base URL. For every domain, we determine what fraction of its web pages were retrieved during the initial collection phase. Domains in which more than 10\% of the web pages have already been gathered are designated as math-related (such as \url{mathoverflow.net}). Next, we conduct manual annotation of URLs within these math-related domains that correspond to mathematical content (e.g., \url{mathoverflow.net/questions}). Any web pages associated with these annotated URLs, which have yet to be collected, are subsequently incorporated into the seed corpus. This procedure supplies a broader set of positive examples, permitting the fastText model to be retrained for higher recall of mathematical material in the subsequent collection cycle. After four such cycles, we ultimately obtain 35.5M mathematical web pages, encompassing 120B tokens. By the fourth round, we observe that nearly 98\% of the web pages had already been acquired during the third round, prompting us to halt additional data collection at this stage.

In order to prevent benchmark data leakage, we adopt the filtering approach of \cite{deepseek-coder}, excluding web pages that include questions or answers from English mathematical evaluation sets such as GSM8K \citep{gsm8k} and MATH \citep{MATH}, as well as Chinese datasets like CMATH \citep{wei2023cmath} and AGIEval \citep{agieval}. Specifically, we remove any portion of text from our math training data that contains an exact match to any 10-gram sequence present in these benchmarks. For benchmark passages with fewer than 10 but at least 3 grams, we conduct an exact string match to identify and eliminate any instances of overlap in the web-sourced material.

\subsection{Validating the Quality of the DeepSeekMath~Corpus}

We run pre-training experiments to investigate how the DeepSeekMath~Corpus is compared with the recently released math-training corpora:

\begin{itemize}[topsep=0pt]
    \item \textbf{MathPile} \citep{mathpile}: a diverse dataset comprising 8.9B tokens, compiled from various resources such as textbooks, Wikipedia, ProofWiki, CommonCrawl, StackExchange, and arXiv—with arXiv accounting for more than 85\% of the material;
    \item \textbf{OpenWebMath} \citep{openwebmath}: a dataset derived from CommonCrawl through targeted filtering for mathematical content, containing a total of 13.6B tokens;
    \item \textbf{Proof-Pile-2} \citep{llemma}: a collection featuring mathematical documents from OpenWebMath, AlgebraicStack (encompassing 10.3B tokens of math-related code), and arXiv articles (28.0B tokens in total). In accordance with \cite{llemma}, experiments conducted on Proof-Pile-2 maintain an arXiv:Web:Code proportion of 2:4:1.
\end{itemize}

\subsubsection{Training Setting}
\label{sec:quality-policy}
We perform specialized mathematical instruction on a pre-existing, general-purpose language model containing 1.3B parameters, utilizing the same architectural design as the DeepSeek LLMs \citep{deepseek-llm}, referred to here as DeepSeek-LLM 1.3B. Distinct models are individually fine-tuned on each specific mathematical dataset, each undergoing 150B tokens of training. The entirety of the training process employs the resource-efficient HAI-LLM \citep{haillm} training framework. Consistent with methodologies adopted for DeepSeek LLMs, optimization is handled via AdamW \citep{adamW} with settings $\beta_1=0.9$, $\beta_2=0.95$, and $\mathrm{weight\_decay}=0.1$. The learning rate follows a staged schedule: it linearly increases during 2,000 warmup steps to its peak, drops to 31.6\% of that peak value after 80\% of the total training steps, and tapers further to 10.0\% by 90\% completion. The highest learning rate applied is capped at 5.3e-4, with training performed using batches of 4M tokens and a context window of 4K tokens.

\subsubsection{Evaluation Results}

\textbf{The DeepSeekMath~Corpus is of high quality, covers multilingual mathematical content, and is the largest in size.}

\begin{itemize}[topsep=0pt]
    \item \textbf{High-quality}:
    We assess downstream capabilities across 8 distinct mathematical evaluation sets employing few-shot chain-of-thought prompting \cite{cot}. Table \ref{tab:corpora_comparison} presents a noticeable superiority for the model trained on the DeepSeekMath~Corpus compared to alternatives. Furthermore, Figure \ref{fig:corpora_comparisons} demonstrates that models trained with DeepSeekMath~Corpus outperform those trained on Proof-Pile-2 after exposure to 50B tokens (equivalent to a complete epoch of Proof-Pile-2), illustrating the higher overall quality of DeepSeekMath~Corpus data. 
    \item \textbf{Multilingual}:
    The DeepSeekMath~Corpus contains resources in a variety of languages, with English and Chinese being the two dominant languages. Table \ref{tab:corpora_comparison} highlights that models trained using DeepSeekMath~Corpus achieve superior mathematical reasoning outcomes in both English and Chinese. By comparison, prior mathematical corpora that focus mostly on English demonstrate limited gains and can even deteriorate results for Chinese-language mathematics reasoning.
    \item \textbf{Large-scale}:
    With its size vastly exceeding that of previously available mathematical datasets, the DeepSeekMath~Corpus supports broader-scale training. As illustrated in Figure \ref{fig:corpora_comparisons}, the DeepSeek-LLM 1.3B model trained on DeepSeekMath~Corpus not only achieves faster progress during learning but also retains its improvements over time. Conversely, the smaller baseline corpora reach their capacity after several passes in training, leading to an early plateau in model performance.
\end{itemize}

\subsection{Training and Evaluating DeepSeekMath-Base 7B}

In this section, we present DeepSeekMath-Base 7B, a foundational model exhibiting notable proficiency in mathematical reasoning. This model is constructed by initializing with DeepSeek-Coder-Base-v1.5 7B \citep{deepseek-coder} and subsequently undergoing training on 500B tokens. The composition of the training dataset is as follows: 56\% originates from the DeepSeekMath Corpus, 4\% from AlgebraicStack, 10\% from arXiv, 20\% from Github code repositories, while the final 10\% is natural language content in English and Chinese sourced from Common Crawl. Training largely adheres to the configuration detailed in Section \ref{sec:quality-policy}, but employs a peak learning rate of 4.2e-4 and a batch size amounting to 10M tokens.

We present an in-depth evaluation of DeepSeekMath-Base 7B's mathematical proficiency, analyzing its capacity to independently generate complete mathematical solutions, utilize tools for solving mathematical tasks, and perform formal theorem proving. Additionally, we offer a broader overview of the base model's abilities, encompassing its competencies in natural language understanding, reasoning, and programming.

\paragraph{Mathematical Problem Solving with Step-by-Step Reasoning}

We assess DeepSeekMath-Base’s capability in mathematical problem solving through the use of few-shot chain-of-thought prompting \citep{cot}, spanning eight benchmarks in both English and Chinese. The selected benchmarks include tests of quantitative reasoning—such as GSM8K \citep{gsm8k}, MATH \citep{MATH}, and CMATH \citep{wei2023cmath}—as well as multiple-choice assessments like MMLU-STEM \citep{mmlu} and Gaokao-MathQA \citep{agieval}. These cover a broad spectrum of mathematical disciplines and range in difficulty from primary to collegiate levels.

Table \ref{tab:base_math_cot} demonstrates that DeepSeekMath-Base 7B achieves the best results among all open-source base models across the eight evaluated benchmarks. This includes comparisons with prominent models such as Mistral 7B \citep{mistral} and Llemma 34B \citep{llemma}—the latter of which has been specifically enhanced using the Proof-Pile-2 dataset \citep{llemma}. Particularly on the MATH benchmark, designed to reflect competition-level mathematical reasoning, DeepSeekMath-Base outperforms its open-source peers by a margin exceeding 10\% in absolute terms. Furthermore, it surpasses Minerva 540B \citep{minerva}, a much larger (77-fold) closed-source model built atop PaLM \citep{palm} and further pre-trained on mathematical materials.

\paragraph{Mathematical Problem Solving with Tool Use} We assess mathematical reasoning assisted by code on the GSM8K and MATH datasets, employing few-shot program-of-thought prompting as described by \citet{pot,pal}. The models are instructed to address each question by generating Python code, taking advantage of libraries like \textit{math} and \textit{sympy} to handle complex calculations. The program’s output is regarded as the solution. As reported in Table \ref{tab:base_math_pot_proof}, DeepSeekMath-Base 7B surpasses the previous leading model, Llemma 34B.

\paragraph{Formal Mathematics}

Automating the construction of formal proofs plays a critical role in verifying the correctness and dependability of mathematical arguments, while also boosting productivity—a topic that has seen growing interest in recent years. In this study, we assess DeepSeekMath-Base 7B on the informal-to-formal proof generation benchmark introduced in \citep{dsp_proof}. The task entails producing a formalized proof given an informal problem statement, its formal equivalent, and an informal proof outline. For our experiments, we utilize the miniF2F dataset \citep{minif2f}, which serves as a standard for testing formal mathematical reasoning at the Olympiad level. We employ a few-shot prompting approach to prompt the model to output a formal Isabelle proof for each item. Consistent with the methodology in \cite{dsp_proof}, our pipeline has the model produce an initial proof sketch, after which we use Sledgehammer \citep{sledgehammer}, a widely used automated reasoning tool, to complete any missing steps. As illustrated in Table \ref{tab:base_math_pot_proof}, DeepSeekMath-Base 7B achieves impressive results on the task of automating proof formalization.

\paragraph{Natural Language Understanding, Reasoning, and Code}

We assess the capabilities of our models in natural language understanding using MMLU \citep{mmlu}, for reasoning through BBH \citep{bbh}, and in programming proficiency via HumanEval \citep{codex} and MBPP \citep{mbpp}. As reported in Table \ref{tab:understanding_reasoning_code}, DeepSeekMath-Base 7B achieves notable improvements over its predecessor, DeepSeek-Coder-Base-v1.5 \citep{deepseek-coder}, on both MMLU and BBH, demonstrating the beneficial influence of mathematical training on language understanding and reasoning tasks. Furthermore, with the inclusion of code tokens in its continued pretraining, DeepSeekMath-Base 7B retains the coding benchmark performance of DeepSeek-Coder-Base-v1.5 on HumanEval and MBPP. In summary, DeepSeekMath-Base 7B delivers considerably better results than the generalist Mistral 7B \citep{mistral} across the three reasoning and programming evaluations.

\section{Supervised Fine-Tuning}

\subsection{SFT Data Curation}

We develop a diverse mathematical instruction-tuning dataset that spans English and Chinese questions from multiple branches of mathematics and encompasses a range of difficulty levels. Each problem is accompanied by a corresponding solution presented in various reasoning formats, including chain-of-thought (CoT) \citep{cot}, program-of-thought (PoT) \citep{pot,pal}, and tool-integrated reasoning format \citep{tora}. The compiled dataset consists of a total of 776,000 training samples.

\begin{itemize}[topsep=0pt]
    \item \textbf{English mathematical datasets}: We provide tool-integrated annotated solutions for both GSM8K and MATH problem sets. Additionally, we include a portion of MathInstruct \citep{MathInstruct} as well as the training split from Lila-OOD \citep{lila}, which feature problems addressed using CoT or PoT approaches. The English dataset encompasses a broad spectrum of mathematical disciplines, such as algebra, probability, number theory, calculus, and geometry.
    \item \textbf{Chinese mathematical datasets}: We assemble a collection of Chinese K-12 math problems, representing 76 distinct sub-fields including, for instance, linear equations. Solutions for these are provided in both CoT and tool-integrated reasoning formats.
\end{itemize}

\subsection{Training and Evaluating DeepSeekMath-Instruct 7B}

This section presents DeepSeekMath-Instruct 7B, developed through mathematical instruction tuning on top of DeepSeekMath-Base. For training, examples are concatenated at random up to a maximum context window of 4K tokens. The model is trained for 500 steps, utilizing a batch size of 256, and employs a fixed learning rate of 5e-5.

We evaluate models' mathematical performance both without and with tool use, on 4 quantitative reasoning benchmarks in English and Chinese. We benchmark our model against the leading models of the time:

\begin{itemize}[topsep=0pt]
    \item \textbf{Closed-source models} comprise: (1) the GPT series, with GPT-4 \citep{gpt4} and GPT-4 Code Interpreter~\footnote{\url{https://openai.com/blog/chatgpt-plugins\#\#code-interpreter}} identified as the most advanced; (2) Gemini Ultra and Pro \citep{gemini}; (3) Inflection-2 \citep{inflection-2}; (4) Grok-1~\footnote{\url{https://x.ai/model-card}}; in addition to several models launched by Chinese firms such as (5) Baichuan-3~\footnote{\url{https://www.baichuan-ai.com}} and (6) the most recent GLM-4~\footnote{\url{https://open.bigmodel.cn/dev/api\#glm-4}}

    from the GLM family \citep{glm}.

These general-purpose models have typically been subjected to a range of alignment processes.  
\item \textbf{Open-source models} span both broad-application and math-specialized categories: general examples are (1) DeepSeek-LLM-Chat 67B \citep{deepseek-llm}, (2) Qwen 72B \citep{qwen}, (3) SeaLLM-v2 7B \citep{seallm}, and (4) ChatGLM3 6B \citep{chatglm3}. Models with targeted improvements in mathematical ability include (5) InternLM2-Math 20B~\footnote{\url{https://github.com/InternLM/InternLM-Math}}, an extension of InternLM2 that receives math-focused instruction tuning after further training, (6) Math-Shepherd-Mistral 7B, created by applying PPO \citep{schulman2017proximal} to Mistral 7B \citep{mistral} with a reward model driven by process-level supervision, (7) the WizardMath collection \citep{wizardmath}, which bolsters Mistral 7B and Llama-2 70B \citep{llama2} mathematical skills using evolve-instruct—a variant of instruction tuning that leverages instructions produced by AI—and PPO, with data drawn mainly from GSM8K and MATH; (8) MetaMath 70B \citep{metamath}, resulting from Llama-2 70B fine-tuning on enriched GSM8K and MATH datasets; (9) ToRA 34B \cite{tora}, based on CodeLlama 34B but further refined for integrated mathematical tool use; and (10) MAmmoTH 70B \citep{MathInstruct}, an instruction-tuned Llama-2 70B leveraging MathInstruct for improved mathematical problem-solving.
\end{itemize}

Table \ref{tab:sft_rl_math} presents the results when tool use is prohibited, indicating that DeepSeekMath-Instruct 7B excels in executing multi-step reasoning. In particular, within the competition-grade MATH benchmark, our model achieves an absolute improvement of no less than 9\% over all open-source contenders, as well as most closed-source systems (such as Inflection-2 and Gemini Pro). This advantage persists even against significantly larger models (for example, Qwen 72B) and those that have undergone targeted reinforcement learning tuned for mathematics (for instance, WizardMath-v1.1 7B). DeepSeekMath-Instruct attains comparable performance to leading Chinese proprietary systems like GLM-4 and Baichuan-3 on the MATH dataset; nonetheless, it falls short of state-of-the-art proprietary models such as GPT-4 and Gemini Ultra.

In the evaluation scenario that permits models to combine natural language reasoning with programmatic tool usage for solving problems, DeepSeekMath-Instruct 7B nearly achieves 60\% accuracy on the MATH dataset, outperforming all previously released open-source models. For the remaining benchmarks, our model demonstrates performance on par with DeepSeek-LLM-Chat 67B, the earlier leading model that is ten times larger.

\subsection{Group Relative Policy Optimization} Following Supervised Fine-Tuning (SFT), reinforcement learning (RL) has demonstrated significant potential in enhancing the mathematical reasoning skills of large language models (LLMs) \citep{wang2023math,wizardmath}. In this section, we present Group Relative Policy Optimization (GRPO), an RL method designed to be both efficient and impactful.

\subsubsection{From PPO to GRPO} Proximal Policy Optimization (PPO) \citep{schulman2017proximal} is an actor-critic RL algorithm that is widely used in the RL fine-tuning stage of LLMs \citep{ouyang2022training}. In particular, it optimizes LLMs by maximizing the following surrogate objective:

\begin{equation}
\footnotesize
    \mathcal{J}_{PPO}(\theta) = \mathbb{E}_{q \sim P(Q),\, o \sim \pi_{\theta_{old}}(O|q)} \left[ \frac{1}{|o|} \sum_{t=1}^{|o|} \min \left( \frac{\pi_\theta(o_{t} | q, o_{<t})}{\pi_{\theta_{old}}(o_{t} | q, o_{<t})} A_{t},\, \text{clip} \left( \frac{\pi_\theta(o_{t} | q, o_{<t})}{\pi_{\theta_{old}}(o_{t} | q, o_{<t})}, 1 - \varepsilon, 1 + \varepsilon \right) A_{t} \right) \right],
\end{equation}

In this context, $\pi_{\theta}$ and $\pi_{\theta_{old}}$ represent the current and previous policy models, respectively, while $q$ and $o$ denote questions and generated outputs sampled from the question dataset and the previous policy $\pi_{\theta_{old}}$. The parameter $\varepsilon$ is a PPO-specific clipping hyper-parameter employed to enhance training robustness. The term $A_t$ stands for the advantage, which is estimated using Generalized Advantage Estimation (GAE) \citep{gae}, leveraging rewards $\{r_{\ge t}\}$ together with a learned value function $V_{\psi}$. Consequently, PPO necessitates concurrent training of a value function along with the policy network. To prevent excessive reward model optimization, the conventional methodology incorporates a token-wise KL divergence penalty—derived from a reference model—into the reward for each token \citep{ouyang2022training}, that is,

\begin{equation}
    r_{t} = r_\varphi(q, o_{\le t}) - \beta \log\left(\frac{\pi_{\theta}(o_{t}|q, o_{<t})}{\pi_{ref}(o_{t}|q, o_{<t})}\right),
\label{eq:PPO-reward}
\end{equation}

In this context, $r_\varphi$ denotes the reward model, while $\pi_{ref}$ represents the reference model—typically corresponding to the original SFT model. The parameter $\beta$ serves as the weighting factor for the KL divergence penalty.

As the value function employed in PPO is typically another model of comparable size as the policy model, it brings a substantial memory and computational burden. Additionally, during RL training, the value function is treated as a baseline in the calculation of the advantage for variance reduction. While in the LLM context, usually only the last token is assigned a reward score by the reward model, which may complicate the training of a value function that is accurate at each token. To address this, as shown in Figure \ref{fig:grpo}, we propose Group Relative Policy Optimization (GRPO), which obviates the need for additional value function approximation as in PPO, and instead uses the average reward of multiple sampled outputs, produced in response to the same question, as the baseline. More specifically, for each question $q$, GRPO samples a group of outputs $\{o_1, o_2, \cdots, o_G\}$  from the old policy  $\pi_{\theta_{old}}$  and then optimizes the policy model by maximizing the following objective:

\begin{equation}
\footnotesize
\begin{split}
    \mathcal{J}_{GRPO}(\theta) &= \mathbb{E}{[q \sim P(Q), \{o_i\}_{i=1}^G \sim \pi_{\theta_{old}}(O|q)]}  \\
    & \frac{1}{G}\sum_{i=1}^G\frac{1}{|o_i|} \sum_{t=1}^{|o_i|} \Bigg( \min \left[ \frac{\pi_\theta(o_{i,t} | q, o_{i,<t})}{\pi_{\theta_{old}}(o_{i,t} | q, o_{i,<t})} \hat{A}_{i,t}, \text{clip} \left( \frac{\pi_\theta(o_{i,t} | q, o_{i,<t})}{\pi_{\theta_{old}}(o_{i,t} | q, o_{i,<t})}, 1 - \varepsilon, 1 + \varepsilon \right)  \hat{A}_{i,t} \right] \\
    &\hspace{5cm} - \beta \mathbb{D}_{KL}\left[\pi_{\theta} || \pi_{ref}\right] \Bigg),
\end{split}
\label{eq:GRPO-obj}
\end{equation}

where $\varepsilon$ and $\beta$ are hyper-parameters, and $\hat{A}_{i,t}$  is the advantage calculated based on relative rewards of the outputs inside each group only, which will be detailed in the following subsections. The group relative way that GRPO leverages to calculate the advantages, aligns well with the comparative nature of rewards models,  as reward models are typically trained on datasets of comparisons between outputs on the same question. Also note that, instead of adding KL penalty in the reward, GRPO regularizes by directly adding the KL divergence between the trained policy and the reference policy to the loss, avoiding complicating the calculation of  $\hat{A}_{i,t}$. And different from the KL penalty term used in (\ref{eq:PPO-reward}), we estimate the KL divergence with the following unbiased estimator \citep{kl_approx}:

\begin{equation}
\small
    \mathbb{D}_{KL}\left[\pi_{\theta} || \pi_{ref}\right] = \frac{\pi_{ref}(o_{i,t}|q,o_{i,<t})}{\pi_{\theta}(o_{i,t}|q,o_{i,<t})} - \log\frac{\pi_{ref}(o_{i,t}|q,o_{i,<t})}{\pi_{\theta}(o_{i,t}|q,o_{i,<t})} - 1,
\end{equation}

which is guaranteed to be positive.

\begin{algorithm}[t]
  \small
  \caption{Iterative Group Relative Policy Optimization}
  \textbf{Input} initial policy model $\pi_{\theta_{\text{init}}}$; reward models $r_\varphi$; set of task prompts $\mathcal{D}$; 
  hyperparameters $\varepsilon$, $\beta$, $\mu$
  \begin{algorithmic}[1]
    \State Initialize policy model: $\pi_\theta \leftarrow \pi_{\theta_{\text{init}}}$
    \For{iteration = 1, \dots, I}
       \State Set reference model: $\pi_{ref} \leftarrow \pi_{\theta}$
      \For{step = 1, \dots, M}
      \State Draw mini-batch $\mathcal{D}_b$ from $\mathcal{D}$
      \State Set previous policy: $\pi_{\theta_{old}} \leftarrow \pi_{\theta}$ 
      \State For each question $q$ in $\mathcal{D}_b$, sample $G$ candidate outputs $\{o_i\}_{i=1}^G \sim \pi_{\theta_{old}} (\cdot \mid q) $
      \State For all generated outputs $o_i$, obtain rewards $\{r_i\}_{i=1}^{G}$ via $r_{\varphi}$ 
      \State Estimate group-relative advantage $\hat{A}_{i,t}$ for every token $t$ in each output $o_i$
      \For{GRPO iteration = 1, \dots, $\mu$}
        \State Perform an update of $\pi_{\theta}$ to increase the GRPO objective as formulated in Equation \ref{eq:GC-GRPO}
      \EndFor
    \EndFor
    \State Refine $r_\varphi$ through ongoing learning with experience replay. 
    \EndFor 
  \end{algorithmic}
  \textbf{Output} $\pi_\theta$
  \label{alg:iter-grpo}
\end{algorithm}

\subsubsection{Outcome Supervision RL with GRPO} Specifically, given a question $q$, a set of candidate outputs $\{o_1, o_2, \cdots, o_G\}$ is generated using the previous policy $\pi_{\theta_{old}}$. Each output is evaluated by a reward model, producing a corresponding reward for each, denoted by $\mathbf{r}=\{r_1, r_2, \cdots, r_G\}$. To standardize the rewards, the mean of the group is subtracted from each reward, and the result is divided by the group's standard deviation. In outcome supervision, the normalized reward is assigned only at the terminal token of each output $o_i$. This normalized reward serves as the advantage estimate $\hat{A}_{i, t}$ for every token in the output, that is, $\hat{A}_{i, t} = \widetilde{r}_i = \frac{r_i- {\rm mean}(\mathbf{r})}{{\rm std}(\mathbf{r})}$. The policy update is then performed by maximizing the objective specified in equation (\ref{eq:GRPO-obj}).

\subsubsection{Process Supervision RL with GRPO} Relying solely on outcome supervision delivers feedback only at the conclusion of each generated output, which often proves inadequate and inefficient for guiding the policy in intricate mathematical reasoning tasks. In line with the methodology in \cite{wang2023math}, we investigate process supervision, wherein feedback is assigned after each reasoning step rather than at the final step. Concretely, for a given problem $q$ and a collection of $G$ generated solutions $\{o_1, o_2, \cdots, o_G\}$, we leverage a process reward model that assigns a reward to every step within these outputs, resulting in reward sets: $\mathbf{R} = \{ \{r_1^{index(1)},\cdots,r_1^{index(K_1)}\}, \cdots,  \{r_G^{index(1)},\cdots,r_G^{index(K_G)}\} \}$, where $index(j)$ denotes the position of the terminal token of the $j$-th reasoning step, and $K_i$ specifies the number of steps within the $i$-th output. To ensure comparability, we normalize the reward values using their mean and standard deviation, as described by $\widetilde{r}_i^{index(j)} = \frac{r_i^{index(j)} - {\rm mean(\mathbf{R})}}{{\rm std(\mathbf{R})}}$.
Subsequently, under process supervision, the advantage attributed to each token corresponds to the aggregate of normalized rewards across all subsequent steps, mathematically expressed as $\hat{A}_{i, t} = \sum_{index(j) \ge t} \widetilde{r}_i^{index(j)}$.
The policy is then refined by maximizing the objective outlined in equation (\ref{eq:GRPO-obj}).

\subsubsection{Iterative RL with GRPO} During reinforcement learning, the reward model initially used may become inadequate for guiding the latest policy model. To address this, we investigate an iterative approach to RL with GRPO. As depicted in Algorithm \ref{alg:iter-grpo}, this iterative GRPO methodology involves creating new training datasets for the reward model by drawing samples from the current policy model. The reward model is then updated iteratively through a replay mechanism, retaining 10\% of data from previous iterations. Afterward, the policy model itself becomes the reference model, and subsequent training of the policy model occurs with the continually updated reward model.

\subsection{Training and Evaluating DeepSeekMath-RL}

Our reinforcement learning procedure utilizes DeepSeekMath-Instruct 7B as the base model. The RL training dataset comprises approximately 144K questions presented in the chain-of-thought format, specifically covering GSM8K and MATH items extracted from the SFT data. We intentionally exclude other SFT questions during RL in order to assess the effects of reinforcement learning on benchmarks where these data are not available in the RL process. The construction of the reward model training set adheres to the approach outlined in \citep{wang2023math}. The initial reward model is trained on DeepSeekMath-Base 7B, employing a learning rate of 2e-5. The policy model's learning rate in the GRPO setup is established as 1e-6, and we utilize a KL penalty coefficient of 0.04. For each prompt, we generate $64$ sampled completions. We set both the maximum sequence length and batch size during training to 1024. The policy model undergoes just one parameter update following each exploration cycle. Evaluation of DeepSeekMath-RL 7B is conducted on the same suite of benchmarks as DeepSeekMath-Instruct 7B. Within this framework, chain-of-thought tasks drawn from GSM8K and MATH are classified as in-domain, whereas all additional benchmarks are treated as out-of-domain tasks.

Table \ref{tab:sft_rl_math} presents a comparative analysis of open- and closed-source models leveraging both chain-of-thought and tool-integrated reasoning techniques on benchmark datasets in English and Chinese. The main observations are as follows: (1) When employing chain-of-thought reasoning, DeepSeekMath-RL 7B achieves accuracy rates of 88.2\% on GSM8K and 51.7\% on MATH. These results exceed those of any open-source model between 7B and 70B parameters and outperform most closed-source counterparts. (2) It is important to note that DeepSeekMath-RL 7B is exclusively fine-tuned on chain-of-thought instruction data derived from GSM8K and MATH, with DeepSeekMath-Instruct 7B as its initialization point. Even though its training set is limited in scope, DeepSeekMath-RL 7B demonstrates superior performance compared to DeepSeekMath-Instruct 7B across all measured metrics, indicating the substantial benefits provided by reinforcement learning.

\section{Discussion}
This section presents the insights gained from both the pre-training and reinforcement learning (RL) experiments.

\subsection{Lessons Learnt in Pre-Training}

We begin by detailing our experiences with pre-training. Except where explicitly stated, our experiments follow the training configurations specified in Section \ref{sec:quality-policy}. Notably, any mention of the DeepSeekMath Corpus throughout this section pertains to the 89B-token dataset obtained during the second round of data collection.

\subsubsection{Code Training Benefits Mathematical Reasoning}

An often-cited but as yet unproven assumption is that instruction in coding enhances reasoning skills. In this work, we provide a partial answer focused on mathematics: code training enhances models’ capacity for mathematical reasoning, regardless of whether external tools are employed.

To study how code training affects mathematical reasoning, we experimented with the following two-stage training and one-stage training settings:

\textbf{Two-Stage Training}

\begin{itemize}[topsep=0pt]
    \item \textbf{Code Training for 400B Tokens $\rightarrow$ Math Training for 150B Tokens}: DeepSeek-LLM 1.3B undergoes an initial phase of 400B code token training, subsequently followed by an additional 150B tokens focused on math.
    \item \textbf{General Training for 400B Tokens $\rightarrow$ Math Training for 150B Tokens}: For comparative purposes, we conduct a control experiment replacing code tokens in the first phase with general tokens drawn from a broad-scale general dataset built by DeepSeek-AI, aiming to assess whether code pretraining confers an advantage over general corpus pretraining with respect to mathematical reasoning capabilities.
\end{itemize}

\textbf{One-Stage Training}

\begin{itemize}[topsep=0pt]
    \item \textbf{Mathematics-focused Training over 150B Tokens}: DeepSeek-LLM 1.3B is trained on a dataset comprising 150 billion mathematics tokens;
    \item \textbf{Joint Training Using a Combination of 400B Code Tokens and 150B Math Tokens}: Sequential training on math data after code data has been found to adversely affect coding capabilities. To address this, we explore whether simultaneously mixing code and math tokens within a unified training stage can not only enhance mathematical reasoning skills but also mitigate catastrophic forgetting of code-related abilities.
\end{itemize}

\paragraph{Results}
Table \ref{tab:code-math} and Table \ref{tab:code-math-general-eval} present downstream task outcomes across various training configurations.

Program-aided mathematical reasoning gains from code training in both the one-stage and two-stage training paradigms. As indicated in Table \ref{tab:code-math}, applying code training independently within the two-stage approach already leads to substantial improvements in solving Python-based tasks from GSM8K and MATH. Subsequent math-focused training during the second stage brings about additional gains. Notably, in the one-stage regime, concurrently presenting code and math tokens not only counteracts the problem of catastrophic forgetting typically observed in the two-stage process, but also results in enhancements for both coding tasks (see Table \ref{tab:code-math-general-eval}) and for program-aided mathematical reasoning capabilities (see Table \ref{tab:code-math}).

Training on code likewise enhances mathematical reasoning skills even in the absence of tool use. When employing a two-stage training approach, beginning with code training brings about noticeable improvements. Additionally, this phase facilitates more effective math training in the following stage, ultimately yielding superior results. Conversely, merging code and math tokens within a single training phase detracts from mathematical reasoning abilities without tool use. One possible explanation is that DeepSeek-LLM 1.3B, given its comparatively small size, may be unable to adequately learn from both code and mathematical content at the same time.

\subsubsection{ArXiv Papers Seem Ineffective in Improving Mathematical Reasoning}

ArXiv papers are commonly included as a component of math pre-training data \citep{minerva,gpt-f,llemma,mathpile}. However, detailed analysis regarding their impact on mathematical reasoning has not been extensively conducted. Perhaps counter-intuitively, according to our experiments, arXiv papers seem ineffective in improving mathematical reasoning. We experiment with models of different sizes, including DeepSeek-LLM 1.3B and DeepSeek-Coder-Base-v1.5 7B \citep{deepseek-coder}, using arXiv corpora that underwent varied processing pipelines:

\begin{itemize}[topsep=0pt]
    \item \textbf{MathPile} \citep{mathpile}: a collection containing 8.9 billion tokens, curated using a set of heuristic cleaning and filtering strategies, with scientific articles from arXiv comprising more than 85\% of the data;
    \item \textbf{ArXiv-RedPajama} \citep{redpajama}: encompasses all arXiv LaTeX documents, stripped of preambles, comments, macro definitions, and reference lists, amounting to 28.0 billion tokens.
\end{itemize}

For our experiments, we independently trained DeepSeek-LLM 1.3B on 150B tokens and DeepSeek-Coder-Base-v1.5 7B on 40B tokens using each arXiv corpus. The results indicate that incorporating arXiv papers does not lead to substantial gains in mathematical reasoning abilities. When models are trained exclusively on arXiv data, they exhibit no significant progress and in some cases, even a decline in performance on a suite of mathematical benchmarks with varying levels of difficulty. These benchmarks encompass quantitative reasoning datasets such as GSM8K and MATH (see Table \ref{tab:arxiv-cot}), multiple-choice assessments like MMLU-STEM (see Table \ref{tab:arxiv-cot}), as well as formal mathematics datasets like miniF2F (see Table \ref{tab:arxiv_atp}).

However, this conclusion has its limitations and should be taken with a grain of salt. We have not yet studied:

\begin{itemize}[topsep=0pt]
    \item How arXiv tokens influence math-centric tasks outside the present study, such as the informalization of theorems, where formal proofs or statements are re-expressed informally;
    \item The consequences of integrating arXiv tokens with additional data sources;
    \item If advantages derived from arXiv papers persist or become more evident as model size increases.
\end{itemize}

Thus, further exploration is required, which we leave for future studies.

\subsection{Insights of Reinforcement Learning}
\subsubsection{Towards to a Unified Paradigm} Here, we introduce a comprehensive framework for examining a range of training strategies, including SFT, RFT, DPO, PPO, and GRPO. We also present experimental investigations to identify key elements influencing this framework. In broad terms, the gradient of any such method with respect to the parameter $\theta$ may be represented as:

\begin{equation}
    \nabla_{\theta}\mathcal{J}_{\textcolor{red}{\mathcal{A}}}(\theta) = \mathbb{E}[\underbrace{(q,o) \sim \textcolor{red}{\mathcal{D}}}_{Data \ Source}]\left( \frac{1}{|o|} \sum_{t=1}^{|o|}  \underbrace{GC_{{\mathcal{A}}}(q, o, t, \textcolor{red}{\pi_{{rf}}})}_{Gradient \ Coefficient}  \nabla_{\theta}\log \pi_{\theta}(o_t | q, o_{<t})\right).
\label{eq:objective}
\end{equation}

There exist three key components: 1) \textit{Data Source $\mathcal{D}$}, which determines the training data; 2) \textit{Reward Function $\pi_{{rf}}$}, which is the source of the training reward signal; 3) \textit{Algorithm $\mathcal{A}$}: which processes the training data and the reward signal to the gradient coefficient $GC$ that determines the magnitude of the penalty or reinforcement for the data. We analyze several representative methods based on such a unified paradigm:

\begin{itemize}
    \item  \textbf{Supervised Fine-tuning (SFT)}: SFT adapts a pretrained model by training it further on a dataset curated by humans specifically for fine-tuning purposes.
    \item \textbf{Rejection Sampling Fine-tuning (RFT)}: RFT extends the fine-tuning of an SFT-pretrained model by selectively training on outputs generated in response to SFT prompts, retaining only those outputs that are verified to be correct.
    \item \textbf{Direct Preference Optimization (DPO)}: DPO further adjusts the SFT model by fine-tuning it on enhanced outputs generated from the SFT model itself, employing a pairwise loss criterion derived from the DPO framework.
    \item \textbf{Online Rejection Sampling Fine-tuning (Online RFT)}: Unlike standard RFT, Online RFT initializes the policy with the SFT model and continuously fine-tunes using outputs generated on-the-fly by the evolving policy model.
    \item \textbf{PPO/GRPO}: PPO/GRPO begins by initializing the policy from the SFT model, then reinforces it through training on samples drawn directly from the contemporaneous outputs of the policy itself.
\end{itemize}

A summary of the elements constituting these approaches is provided in Table \ref{tab:data_policy}. For a comprehensive explanation of the derivation process, consult Appendix \ref{app:analysis-rl}.

\paragraph{Observation about Data Source} We classify the origin of the data into two types: online sampling and offline sampling. In the case of online sampling, the training data are obtained from the outputs generated by the current policy model during real-time exploration. Conversely, offline sampling refers to collecting training data from the behavior of an initial SFT model. Methods such as RFT and DPO utilize offline sampling, whereas Online RFT and GRPO adopt the online approach.

As illustrated in Figure \ref{fig:rl-analysis}, our results indicate that Online RFT yields markedly better performance than RFT across both benchmarks. In particular, during the initial phases of training, the performance of Online RFT is on par with RFT; however, as training advances, Online RFT secures a clear and increasing lead, underscoring the benefits of online training methods. This pattern is expected: early in training, the actor and SFT model remain very similar, resulting in sampled data with minimal divergence. As training progresses, the actor's outputs begin to diverge more substantially, making real-time data sampling much more advantageous in capturing these differences.

\paragraph{Observation about Gradient Coefficient} The algorithm utilizes the reward signal by converting it into a gradient coefficient, which is then used for updating model parameters. In our studies, the reward function is decomposed into two categories: `Rule' and `Model.' The `Rule' category assesses response quality based on answer accuracy, whereas `Model' involves training a reward model to evaluate and assign scores to each response. Notably, the training instances for the reward model are sourced from the `Rule' assessments. Equations \ref{eq:GC-RFT} and \ref{eq:GC-GRPO} elucidate a central distinction between GRPO and Online RFT: GRPO is characterized by its dynamic adjustment of the gradient coefficient in accordance with the reward values assigned by the reward model. This mechanism enables the model to provide varied reinforcement and penalties aligned with the scale of the received rewards. Conversely, Online RFT is devoid of this capability; it neither imposes penalties for incorrect answers nor varies the reinforcement for correct ones, thereby maintaining uniform intensity of positive reinforcement across all correct responses.

Figure \ref{fig:rl-analysis} illustrates that GRPO outperforms online RFT, underscoring the effectiveness of modulating both positive and negative gradient coefficients. Moreover, GRPO+PS achieves better results than GRPO+OS, demonstrating the advantage of adopting finely-tuned, step-specific gradient coefficients. We also investigate the use of iterative RL by implementing two iterative cycles in our experiments. As depicted in Figure \ref{fig:iter-rl}, the results reveal that iterative RL markedly enhances performance, with the most notable improvement occurring during the initial iteration.

\subsubsection{Why RL Works?} In this study, we apply reinforcement learning to a selected subset of instruction tuning data, observing a notable improvement over models trained solely with instruction tuning. To shed light on the mechanisms underlying the effectiveness of reinforcement learning, we assess both Pass@K and Maj@K accuracy metrics for the Instruct and RL models using two benchmark datasets. Figure \ref{fig:maj-pass} illustrates that reinforcement learning increases Maj@K accuracy, whereas Pass@K remains largely unchanged. These results suggest that reinforcement learning bolsters the robustness of the output distribution, primarily by elevating the likelihood that the correct answer appears within the top $K$ outputs, rather than by fundamentally enhancing the base capabilities of the model. Likewise, \citep{wang2023making} report a \textbf{misalignment problem} related to reasoning in SFT models, and demonstrate that reasoning skills can be improved in these models by applying various preference alignment techniques \citep{yuan2023rrhf,song2023preference,wang2023making}.

\subsubsection{How to Achieve More Effective RL?}
\label{sec:effec_rl}
We show that RL performs quite effectively on tasks involving mathematical reasoning. Furthermore, we introduce a comprehensive framework for interpreting various prominent training approaches. This framework characterizes these methods as variations of direct or reduced RL strategies. As outlined in Equation \ref{eq:objective}, the main elements involved are the Data Source, Algorithm, and Reward Function. Below, we discuss several avenues for future research pertaining to each of these components.

\paragraph{Data Source} The data source constitutes the foundational input for all forms of training approaches. Within the scope of RL, this specifically denotes unlabeled questions, where the policy model generates associated outputs. In this work, our focus is limited to utilizing questions originating from the instruction tuning phase, alongside basic nucleus sampling for output generation. This design choice may contribute to the observation that our RL pipeline yields improvements primarily for Maj@K metrics. Looking ahead, we aim to apply our RL framework to prompts outside the original distribution and incorporate \textbf{advanced sampling (decoding) strategies}, such as those employing tree-search techniques \citep{yao2023tree}. Additionally, the importance of \textbf{efficient inference techniques} \citep{xia-etal-2023-speculative,leviathan2023fast,kwon2023efficient,xia2024unlocking}, which govern how effectively policy models explore the data space, cannot be overstated.

\paragraph{Algorithms} Algorithms utilize both data and the reward signal to determine the gradient coefficient, thereby adjusting model parameters. As described by Equation \ref{eq:objective}, contemporary techniques predominantly \textbf{TRUST} the feedback provided by the reward function to modify the conditional probability of specific tokens, either raising or lowering it. Nevertheless, it is not feasible to guarantee the consistent reliability of the reward signal, particularly when faced with highly complex tasks. For instance, the PRM800K dataset \citep{lightman2023let}, even with meticulous annotations by skilled annotators, still exhibits an error rate of about 20\% in its labeling\footnote{\url{https://github.com/openai/prm800k/issues/12\#issuecomment-1728491852}}. In response, we aim to investigate reinforcement learning algorithms that maintain robustness when confronted with noisy reward signals. We posit that adopting \textbf{WEAK-TO-STRONG} \citep{burns2023weak} alignment methods has the potential to fundamentally transform existing learning paradigms.

\paragraph{Reward Function} The reward function serves as the primary training signal in reinforcement learning (RL), typically realized as a neural reward model. We identify three key research avenues for reward models: 1) \textbf{Strategies to improve reward model generalization.} To truly advance language model capabilities, the reward model must successfully generalize to manage out-of-distribution questions and sophisticated decoding outputs; without this, reinforcement learning risks merely entrenching the existing LLM output distribution. 2) \textbf{Methods to characterize reward model uncertainty.} Accurately quantifying uncertainty could form a critical connection between limited reward models and weak-to-strong learning paradigms. 3) \textbf{Techniques for constructing high-quality process reward models efficiently}, enabling delivery of detailed, fine-grained signals during the reasoning process \citep{lightman2023let,wang2023math}.

\section{Conclusion, Limitation, and Future Work}
We introduce DeepSeekMath, a model that surpasses all open-source counterparts on the challenging MATH benchmark and nearly matches the results of proprietary systems. Built upon DeepSeek-Coder-v1.5 7B, DeepSeekMath is continuously trained with 500B tokens, notably including 120B mathematical tokens primarily extracted from Common Crawl. Our comprehensive ablation investigations reveal that web-sourced content holds considerable promise for acquiring robust mathematical datasets, whereas arXiv data may not yield the anticipated benefits. We propose Group Relative Policy Optimization (GRPO), a modified form of Proximal Policy Optimization (PPO), which delivers substantial improvements in mathematical reasoning while using less memory. Experimental outcomes indicate that GRPO maintains its effectiveness even as DeepSeekMath-Instruct 7B achieves strong benchmark results. Furthermore, we present a unified framework to interpret a suite of methods and outline prospective avenues for more efficient reinforcement learning.

While DeepSeekMath~demonstrates strong results on benchmarks for quantitative reasoning, its performance in geometry and theorem-proving remains inferior compared to proprietary models. For example, our preliminary testing revealed that the model fails to solve questions involving triangles and ellipses, possibly reflecting a pre-training and fine-tuning data selection bias. Moreover, due to its model size constraints, DeepSeekMath~lags behind GPT-4 in few-shot learning abilities; whereas GPT-4's performance benefits from few-shot examples, DeepSeekMath~shows little distinction between zero-shot and few-shot scenarios. Looking ahead, we aim to enhance our data selection mechanisms to assemble a more robust and higher-quality pre-training dataset. Furthermore, we intend to investigate new avenues (Section \ref{sec:effec_rl}) to achieve more effective reinforcement learning in LLMs.

\end{document}